%% ma: L11-B14-0029
%% lop: 11
%% chuong: 4
%% bai: 14
%% dang: 11.14.4
%% loai: TLN
%% muc: VD
%% dap_an: "1"
%% nguon: "tai-lieu-hoc-tap-mon-toan-lop-11-theo-chuong-trinh-giao-duc-pho-thong-2018.pdf (Lớp Toán Thầy Duy, Phạm Lê Duy, THPT Chu Văn An), Chương 4 Bài 5, Phần 3, Câu 5 (tr.535)"
%% kien_thuc: "Lớp 11 Bài 14 (phép chiếu song song)"
%% trang_thai: cho-duyet
%% ly_do: "Dạng mới đề xuất (11.14.x), chờ giáo viên duyệt dạng; tài liệu gốc không có lời giải mẫu, đáp án do kho tự giải và kiểm bằng sympy."
\begin{ex}
Cho hình chóp $S.ABCD$ đáy là hình vuông tâm $O$ cạnh bằng $2$. Trên cạnh $SB,SD$ lần lượt lấy điểm $M,N$ sao cho $SM=2MB$ và $SN=\dfrac13SD$. Diện tích hình chiếu song song của tam giác $AMN$ lên mặt phẳng $(ABCD)$ theo phương của đường thẳng $SO$ bằng
\shortans{1}
\loigiai{\sloppy Theo phương $SO$: $A\mapsto A$, $M\mapsto P\in OB$ với $OP=\dfrac23OB$, $N\mapsto Q\in OD$ với $OQ=\dfrac13OD$ ($P,Q$ nằm hai phía của $O$). Hình vuông cạnh $2$ có $OA=OB=OD=\sqrt2$. $PQ=OP+OQ=\dfrac{2\sqrt2}3+\dfrac{\sqrt2}3=\sqrt2$ và $AO\perp PQ$ nên diện tích $\triangle APQ=\dfrac12\cdot AO\cdot PQ=\dfrac12\cdot\sqrt2\cdot\sqrt2=1$.}
\end{ex}
